\section{Conclusion}

Using an earthquake catalogue from the USGS database, we compared a
homogeneous Poisson process with several Hawkes process specifications,
including marked and unmarked exponential and power-law kernels. Simple
diagnostics on the raw inter-event times reveal pronounced temporal
clustering that is incompatible with a constant-rate Poisson model.
Residual analyses under the time-rescaling theorem, based on
Kolmogorov--Smirnov and chi-squared tests, point to the marked
exponential Hawkes process as the most adequate representation of the
data: it captures both short-term aftershock clustering and the positive
dependence of triggering intensity on earthquake magnitude, yielding
rescaled inter-arrivals that are approximately $\mathrm{Exp}(1)$.

Several limitations remain. The analysis is purely temporal and ignores
the spatial distribution of events; a space--time Hawkes process would
likely provide a more faithful description of seismicity in the region.
The background rate $\mu$ is assumed constant throughout the observation
window, which may not hold over long periods or across distinct seismic
sequences.

